\section{Conclusion}
\label{sec:conclusion}

State legislative redistricting is the most consequential and least studied
domain of algorithmic redistricting. The F research track provides the first
systematic application of minimum-edge-cut recursive bisection to all 50 state
legislatures, extending the \texttt{redist} pipeline with chamber-aware
configuration and resolution-adaptive geographic units.

The central methodological contribution of the F track is the resolution
selection rule: use block-group resolution when $k/n_{\mathrm{tracts}} > 0.05$.
This simple threshold, derived from the requirement that each district contain
at least 20 base units on average, correctly classifies all 50 state chambers
and enables the pipeline to handle district counts from $k=20$ (Alaska Senate)
to $k=400$ (New Hampshire House) without modification.

The central empirical finding is that algorithmically generated state house maps
are \emph{more compact} than their congressional counterparts by the Polsby-Popper
measure ($\overline{PP} = 0.401$ vs.\ $0.361$), while achieving tighter
population balance ($\leq 0.5\%$) than the legally required standard
($\leq 5\%$) and splitting fewer counties than enacted plans.

These findings establish that the same algorithmic principles that produce
compact, population-balanced congressional maps---minimum-edge-cut recursive
bisection without political inputs---scale gracefully to the more demanding
state legislative context. The F track lays the foundation for a broader
research program in algorithmic state governance: once a state's geographic
and demographic data are loaded into the \texttt{redist} pipeline, generating
a principled, defensible legislative map is a matter of running
\texttt{redist state --chamber house --resolution block\_group}.

\paragraph{Data and replication.}
All results in the F track are generated by the \texttt{redist} Rust binary
(version tagged with the paper's release). Block-group adjacency data for
the states studied in F.1--F.3 are available at the project repository.
Full 50-state block-group adjacency data require approximately 4~GB of storage
and can be generated using the \texttt{build\_unit\_adjacency.py} script.
